\section*{v2.9.0 Spectral Readout Filters and Trace-Compatible Mode Selection Core}
\addcontentsline{toc}{section}{v2.9.0 Spectral Readout Filters and Trace-Compatible Mode Selection Core}

\begin{quote}
\textbf{Normalization note (v2.17.1 local review).}
This block is retained as a \emph{historical spectral-readout and trace-filter precursor}. It contributes proof memory and specialized readout discipline, but should be read through the later normalized readout and interface phases.
\end{quote}

The spectral admissibility layer from v2.8.0 decides which generators may enter the
programme at all. The next structural question is subtler: once a generator is
admissible, which parts of its spectrum are actually readable in a way compatible
with the ledger, HC confinement, probe comparison, and later trace-style summary
principles? In this version the monolith therefore adds a spectral readout layer.
Not every admissible mode is automatically promoted to a structural observable. A
mode must also survive a filter test that distinguishes stable, tagged, and
non-readable contributions.

\subsection*{1. Readout filter principle}
Let $\mathcal G$ be spectrally admissible on the HC-confined state space
$\mathfrak X$, and let $\Sigma_{\mathrm{adm}}(\mathcal G)$ denote the admissible part
of the spectrum. We introduce a family of bounded readout filters
$\{\mathcal F_\alpha\}_{\alpha\in A}$ acting on the admissible sector, where each
$\mathcal F_\alpha$ is designed to isolate one structural mode class: triadic,
orthogonal, measurement, residual, or probe-induced. A mode contribution is called
\emph{trace-compatible} if there exists at least one filter $\mathcal F_\alpha$ such
that
\[
\mathcal F_\alpha \Pi_{\mathrm{HC}} u = \Pi_{\mathrm{HC}} u + r_\alpha,
\qquad
\|r_\alpha\|_{\mathfrak X} \le \varepsilon_\alpha \|u\|_{\mathfrak X},
\]
with either small remainder $r_\alpha$ or explicitly tagged remainder. Thus the
readout is not merely spectral isolation but structural readability modulo tagged
error.

\paragraph{Principle 3 (filtered observability).}
A spectral component may be used in the proof architecture only if it is both
spectrally admissible and trace-compatible under an approved readout filter family.
This prevents uncontrolled spectral clutter from being mistaken for genuine
structural information.

\subsection*{2. Trace-compatible decomposition}
The admissible spectrum should now admit a decomposition of the schematic form
\[
\Sigma_{\mathrm{adm}}(\mathcal G)
= \Sigma_{\mathrm{read}} \cup \Sigma_{\mathrm{tag}} \cup \Sigma_{\mathrm{mute}},
\]
where readable modes lie in $\Sigma_{\mathrm{read}}$, tagged but non-primary modes lie
in $\Sigma_{\mathrm{tag}}$, and the remainder $\Sigma_{\mathrm{mute}}$ is admissible but
not allowed to enter the principal structural trace. The point is not that mute
modes vanish, but that they cannot serve as primary carriers of theorem-level
readout. This sharpens the difference between existence and interpretability.

\subsection*{3. HC-filter compatibility}
The HC projector from v2.8.0 is now required to commute with readout at least up to
controlled tagged remainder. Concretely, one seeks estimates of the form
\[
\mathcal F_\alpha \Pi_{\mathrm{HC}}
= \Pi_{\mathrm{HC}} \mathcal F_\alpha + \mathcal T_\alpha,
\]
where $\mathcal T_\alpha$ is either small in the admissible norm or belongs to a
source-tagged correction class. This means HC confinement is preserved under mode
selection: filtering is not allowed to re-open forbidden channels except through a
ledger-visible correction.

\subsection*{4. Mode budgets and readable trace mass}
The budget language from v2.5.0 now acquires a spectral refinement. For every
approved filter $\mathcal F_\alpha$ define the readable trace mass
\[
\mathfrak M_\alpha(u) := \mathcal E(\mathcal F_\alpha u)
\]
with $\mathcal E$ the field-energy functional from v2.7.0. Structural readability
requires a balance inequality of the schematic form
\[
\sum_{\alpha\in A_{\mathrm{read}}} \mathfrak M_\alpha(u)
\le C_{\mathrm{tr}}
\Bigl( \mathcal E(u) + \mathrm{Tag}(u) + \mathrm{Leak}(u) \Bigr).
\]
Hence readable spectral mass cannot exceed the ledger budget except through already
visible tagged or leakage terms. This gives a quantitative bridge between spectral
filtering and earlier channel transport estimates.

\subsection*{5. Probe-sector mode inheritance}
For the two-probe subsystem the same distinction becomes a comparison principle.
Let $\mathcal G_{\mathrm{probe}}$ be the effective probe generator and
$\mathcal F^{\mathrm{probe}}_\beta$ its admissible filter family. Then every readable
probe mode should either descend from a readable ambient mode or be marked as a
probe-only tagged correction. In particular, one asks for comparison bounds of the
form
\[
\mathfrak M^{\mathrm{probe}}_\beta(v)
\le C_{\mathrm{cmp}}
\bigl( \mathfrak M_\alpha(u) + \mathrm{Mismatch}(u,v) \bigr),
\]
for matched ambient/probe states $(u,v)$. This prevents the probe sector from
manufacturing fake observables that have no ambient structural parent.

\subsection*{6. Consequence for future field equations and appendix tools}
The admission checklist for candidate field equations is now extended once more:
\begin{enumerate}[label=(\arabic*)]
  \item verify semigroup dissipativity and ledger closure;
  \item verify spectral admissibility and resolvent compatibility;
  \item construct an HC-compatible readout filter family;
  \item separate readable, tagged, and mute spectral sectors; and
  \item prove probe inheritance of readable modes.
\end{enumerate}
The appendix tool taxonomy is correspondingly refined. Besides resolvent and gap
certificates, a strong tool may now also be labelled as readout filter constructor,
trace-mass comparator, mute-sector exclusion lemma, or probe-inheritance witness.
These labels indicate that the tool produces not only existence information but also
structural legibility.

\subsection*{7. v\docversion\ readiness criterion}
The present step should be regarded as successful if the following can now be done:
\begin{itemize}
  \item admissible spectra can be split into readable, tagged, and mute parts;
  \item HC confinement survives mode filtering modulo controlled tagged remainder;
  \item readable trace mass is bounded in the same ledger language as channel budgets;
  \item probe readout inherits ambient readable modes rather than inventing new ones; and
  \item future field equations can be screened not only for existence, but for structural legibility.
\end{itemize}

This is the gain of v\docversion. The companion monolith now carries not only
closure, channel specialization, transport dynamics, forward invariance,
field-energy realization, semigroup readout, and spectral admissibility, but also a
first filtered-observability layer. Future field equations are therefore constrained
three times: they must generate the correct ledger-dissipative semigroup, they must
belong to the correct resolvent class, and their admissible modes must pass a
trace-compatible readout filter before they can count as structural content.


